\section{语音识别}
\label{sec:nlp-speech-recognition}

\term{自动语音识别}（automatic speech recognition, ASR）把一段声音转换为规定形式的文字。沿用第\ref{sec:nlp-speech-representation}节的记号：声学特征矩阵$\bm{X}$经编码器得到$T$个编码步的表示$\bm{h}_1,\ldots,\bm{h}_T$，目标文本为$y=(y_1,\ldots,y_U)$，符号来自词表$V$，可以是字符、子词或词。识别器要给出条件概率$p(y\mid\bm{X})$，并在推断时找出概率高的文本。难点在于训练数据只给出整段声音与整句文字，没有说明哪个字对应哪几帧。本节各方法的根本差别，就在于怎样处理这种未给定的对应关系。

\subsection{识别的挑战与模型谱系}
\label{sec:nlp-speech-asr-challenges}

一个识别样本至少包含声音、采样率、语种和转写规范。转写规范规定输出单位，以及数字、语气词、标点和大小写如何书写；这些规定不改变声音，却决定“正确答案”是什么。语音识别面临的主要挑战可以归为五项。

\begin{enumerate}
  \item 长度不等且没有帧级对齐。第\ref{sec:nlp-speech-features}小节的例子中，一秒声音有98个特征帧，降采样后约25个编码步，对应的中文可能只有四五个字。语速改变时，同一个字占用的帧数随之改变。训练标签通常只有整句文字，模型必须自己决定对应关系，或者把所有可能的对应都考虑进去。
  \item 上下文与流式的矛盾。听到“四月十”时还不能确定后面是“七日”还是“八日”，同音词也常要靠后文区分，所以看见右侧声音能提高准确率。在线系统却必须在用户说完之前陆续交出文字，可等待的右侧声音受延迟预算限制。
  \item 配对标注稀缺。人工转写一小时录音往往需要数倍于录音时长的工作量。网络字幕数量大，但含有错位、删节、意译和非语音文字，属于弱监督。
  \item 噪声、口音与设备差异。背景噪声、混响、口音和麦克风改变了观测分布而不改变文字，在一种条件下训练的模型换到另一种条件时错误率会上升。
  \item 多说话者。会议录音中多人交替甚至同时说话，系统还要回答“谁在何时说话”。这一项在第\ref{sec:nlp-speech-xvector}小节单独处理。
\end{enumerate}

评价识别结果先要比较两个符号序列。设参考序列为$r_1,\ldots,r_N$，识别结果为$\hat r_1,\ldots,\hat r_M$，$d_{i,j}$表示把参考前$i$项变成识别结果前$j$项所需的最少编辑次数，则
\begin{equation}
 d_{i,j}=\min\{d_{i-1,j}+1,\ d_{i,j-1}+1,\
              d_{i-1,j-1}+\mathbf{1}[r_i\ne\hat r_j]\},
 \quad d_{i,0}=i,\ d_{0,j}=j.
 \label{eq:nlp-speech-edit-distance}
\end{equation}
三项依次对应删除、插入与替换或匹配。回溯最短路径得到替换数$N_{\mathrm{S}}$、删除数$N_{\mathrm{D}}$和插入数$N_{\mathrm{I}}$，\term{词错误率}（word error rate, WER）定义为$(N_{\mathrm{S}}+N_{\mathrm{D}}+N_{\mathrm{I}})/N$。中文通常按字计算，称为\term{字符错误率}（character error rate, CER）。插入可以使错误率超过1；参考为空时只能报告插入数。

计分前必须固定规范化口径。参考写“14:00”而识别结果写“下午两点”，按字比较会产生多次编辑，语义却相同；繁简、全半角和标点是否计分也要事先约定。另一方面，平均错误率会掩盖关键字段的失败：参考“四月十八日”被识别为“四月十七日”，CER为1/5；放在一条五十字的通知中只占2\%，通知的核心信息却已经错了。因此，日期、数量和专名这类字段宜在CER之外单独统计。

图\ref{fig:nlp-speech-asr-lineage}给出本节的方法谱系。HMM路线用显式的声学状态描述“在某个音上停留，再前进到下一个音”，需要帧级对齐和多阶段训练。CTC在所有单调对齐上求和，只需整句标签，但假设各步输出条件独立。注意力编码器—解码器读取已输出的前缀，去掉了条件独立假设，却失去了单调性保证，也难以流式运行。联合CTC与注意力训练用CTC分支找回单调约束；RNN-T则在同一个格点上同时保留路径求和与输出历史。Whisper不改变基本结构，而是用大规模弱监督数据和任务词元扩大覆盖。流式识别最后限制编码器可见的声音范围。第\ref{sec:nlp-speech-asr-comparison}小节把这些差别汇总为对照表。

\begin{figure*}[htbp]
  \centering
  % 第九卷 15.2.1 节：语音识别模型谱系（跨栏，约156 mm）。
% 蓝色：显式状态对齐；黄色：对齐路径求和；红色：注意力条件生成。
% 实线箭头表示“改变了前一方法的哪个机制”，虚线表示另一来源的组成部分。
\begin{tikzpicture}[
  x=1mm,y=1mm,
  figure picture,
  base/.style={figure node,minimum width=30mm,minimum height=11mm,
    inner ysep=2pt,font=\figurefont\scriptsize},
  nst/.style={base,draw=FigureInput,fill=FigureInputFill},
  nps/.style={base,draw=FigureModel,fill=FigureCream},
  nat/.style={base,draw=FigureOutput,fill=FigureOutputFill},
  fl/.style={figure flow},
  dl/.style={figure flow,draw=FigureMuted,dash pattern=on 2pt off 1.5pt},
  el/.style={font=\figurefont\scriptsize,text=FigureMuted,inner sep=1pt,align=center}
]
  \node[nst] (hmm)   at (15,40) {HMM--GMM\\帧级状态与GMM似然};
  \node[nst] (hyb)   at (15,18) {混合HMM--DNN\\网络给出状态后验};
  \node[nps] (ctc)   at (57,40) {CTC\\对单调路径求和};
  \node[nps] (rnnt)  at (57,18) {RNN-T\\路径求和＋输出历史};
  \node[nps] (chunk) at (57,-4) {分块流式编码器\\限制右侧可见范围};
  \node[nat] (aed)   at (99,40) {注意力编码器—解码器\\读取历史，无单调约束};
  \node[nat] (joint) at (99,18) {联合CTC/注意力\\CTC分支补单调约束};
  \node[nat] (wsp)   at (141,40) {Whisper\\弱监督与任务词元};
  \node[nat] (wss)   at (141,18) {Whisper-Streaming\\稳定前缀提交};

  \draw[fl] (hmm) -- node[el,left] {神经网络\\替换GMM} (hyb);
  \draw[fl] (hyb.north east) -- node[el,below right=-1mm and 0mm,pos=.45] {去掉帧级\\对齐标签} (ctc.south west);
  \draw[fl] (ctc) -- node[el,above] {引入\\输出历史} (aed);
  \draw[fl] (ctc) -- node[el,right] {加入\\预测网络} (rnnt);
  \draw[dl] (aed.south west) -- (rnnt.north east);
  \draw[fl] (aed) -- node[el,right] {共享编码器\\加CTC分支} (joint);
  \draw[fl] (aed) -- node[el,above] {扩大\\数据覆盖} (wsp);
  \draw[fl] (wsp) -- node[el,right] {缓冲重识别} (wss);
  \draw[fl] (rnnt) -- node[el,right] {因果化\\与分块} (chunk);

  % 图例
  \node[nst,minimum width=6mm,minimum height=4mm] (l1) at (88,-4) {};
  \node[anchor=west,font=\figurefont\scriptsize] at (l1.east) {显式状态对齐};
  \node[nps,minimum width=6mm,minimum height=4mm] (l2) at (111,-4) {};
  \node[anchor=west,font=\figurefont\scriptsize] at (l2.east) {对齐路径求和};
  \node[nat,minimum width=6mm,minimum height=4mm] (l3) at (134,-4) {};
  \node[anchor=west,font=\figurefont\scriptsize] at (l3.east) {注意力条件生成};
\end{tikzpicture}

  \caption{语音识别模型谱系。颜色表示处理声文对应的方式：蓝色用显式状态对齐，黄色对所有单调对齐路径求和，红色用注意力在生成时读取声音。实线箭头旁注明后一方法改变了前一方法的哪个机制；虚线表示RNN-T同时吸收了注意力模型读取输出历史的做法。CTC同样可以接分块流式编码器，图中只画出RNN-T一支。}
  \label{fig:nlp-speech-asr-lineage}
\end{figure*}

\subsection{HMM识别路线}
\label{sec:nlp-speech-hmm}

最早系统处理变长对应的思路，是让每个发音单位在若干声学状态中停留若干帧，再按顺序前进到下一个状态。设$W$为词序列，识别要找后验概率最大的词序列。由贝叶斯公式，$p(\bm{X})$对所有候选相同，可以略去：
\begin{equation}
 \hat W=\argmax_W p(W\mid\bm{X})
       =\argmax_W p(\bm{X}\mid W)\,P(W).
 \label{eq:nlp-speech-noisy-channel}
\end{equation}
$p(\bm{X}\mid W)$由\term{声学模型}（acoustic model）和\term{发音词典}（pronunciation lexicon）共同给出：词典把词展开成音素串，每个音素对应一个小HMM，声学模型评价每帧特征在状态下的似然。$P(W)$由第\ref{sec:nlp-classical-ngram-models}节的$n$元语言模型给出。三者最终组合成一张统一的解码图，在图上搜索最优路径。

HMM的前向算法、Viterbi算法和Baum--Welch估计已在第\ref{sec:nlp-classical-hmm}小节以词性标注为例推导过；带缩放的前向后向算法和期望计数的一般推导可参考\BookChapterRef{chap:pgm-dynamic}{卷六动态有向模型章}。本小节只写识别特有的部分：音素状态的拓扑、高斯混合似然、上下文相关状态的绑定，以及解码图的组合。本小节以$t$索引特征帧（帧移10\,ms，不降采样），$T$为帧数。

\subsubsection{HMM--GMM声学模型}
\label{sec:nlp-speech-hmm-gmm}

每个音素用一个三状态从左到右的HMM表示，三个状态大致对应音的起始、稳定段与过渡段。状态$j$只允许自环或前进到$j+1$，转移概率为$a_{j,j}$和$a_{j,j+1}=1-a_{j,j}$。自环让一个状态持续任意帧：停留恰好$k$帧的概率为$a_{j,j}^{k-1}(1-a_{j,j})$，平均停留$1/(1-a_{j,j})$帧。例如$a_{j,j}=0.8$时平均停留5帧，即50\,ms。这种几何分布的时长模型很粗糙，但它让同一个HMM能适应不同语速。按词典把音素HMM首尾相接，就得到词的HMM；再按词序列相接，词间可插入可选的静音模型，就得到整句的HMM。

状态$j$下观测特征$\bm{x}_t\in\mathbb{R}^{d}$的似然用\term{高斯混合模型}（Gaussian mixture model, GMM）表示。传统系统常取$d=39$，即13维MFCC及其一阶、二阶差分。设状态$j$有$M$个分量，权重$c_{j,m}\geq0$、$\sum_mc_{j,m}=1$，均值$\bm{\mu}_{j,m}$，对角协方差$\operatorname{diag}(\bm{\sigma}_{j,m}^2)$，则
\begin{equation}
\begin{aligned}
 b_j(\bm{x}_t)&=\sum_{m=1}^{M}c_{j,m}\,
   \mathcal{N}\bigl(\bm{x}_t;\bm{\mu}_{j,m},\operatorname{diag}(\bm{\sigma}_{j,m}^2)\bigr),\\
 \log\mathcal{N}(\cdot)&=-\frac12\sum_{k=1}^{d}
   \left[\log(2\pi\sigma_{j,m,k}^2)+\frac{(x_{t,k}-\mu_{j,m,k})^2}{\sigma_{j,m,k}^2}\right].
\end{aligned}
\label{eq:nlp-speech-gmm}
\end{equation}
对角协方差把每个分量的参数从$O(d^2)$降到$O(d)$，计算也只需逐维累加。这一简化依赖特征各维近似去相关；第\ref{sec:nlp-speech-features}小节MFCC中的离散余弦变换正起这个作用，这也是MFCC长期与GMM搭配的原因。

识别时在整句HMM上运行Viterbi算法。用对数分数表示，到第$t$帧处于状态$s$的最优路径分数为
\begin{equation}
 v_t(s)=\log b_s(\bm{x}_t)+
       \max_{r\to s}\{v_{t-1}(r)+\log a_{r,s}\},
 \label{eq:nlp-speech-hmm-viterbi}
\end{equation}
$r\to s$取遍允许转移到$s$的前驱。起始状态取对数概率0，其余取$-\infty$；保存取最大值的前驱，在允许的终态中选最佳分数并回溯。词边界处再加上语言模型的对数分数。若教学词典把“hi”映射成两个音素/h/、/ay/，一条可能的回溯路径是$h_1,h_1,h_2,h_3,ay_1,ay_2,ay_2,ay_3$；沿词典读出的是一个词“hi”，八个状态并不对应八个输出符号。

同一个音素在不同邻音之间的发音差别很大，例如/ay/在/h/之后与在/m/之后的起始段不同。\term{三音素}（triphone）模型为每个“左邻—中心—右邻”组合单独建模。设音素集有40个，三音素有$40^3=64000$种，每种3个状态，共约19万个状态；若每个状态用16个39维对角高斯，参数约2.4亿，且大多数组合在训练数据中从未出现。\term{状态绑定}（state tying）用决策树解决这个问题：对同一中心音素同一位置的所有三音素状态，从根节点开始，用“左邻是否为鼻音”“右邻是否为元音”这类语音学问题逐层二分，每次选择使训练数据对数似然增益最大的问题，增益低于阈值或节点数据量过少时停止\scite{nlp-speech-young1994}。叶节点就是绑定后的状态，通常只有数千个；训练中没见过的三音素沿树的问题走到某个叶节点，也能得到声学模型。

参数用Baum--Welch算法估计。每条训练语句按转写和词典拼成整句HMM，路径只在这张图中展开，这种做法称为嵌入式训练。E步用前向后向算法算出状态占用$\gamma_t(j)$，再按分量似然分配到各个分量：$\gamma_t(j,m)=\gamma_t(j)\,c_{j,m}\mathcal{N}_{j,m}(\bm{x}_t)/b_j(\bm{x}_t)$。M步用加权统计更新参数：
\begin{equation}
\begin{aligned}
 \bm{\mu}_{j,m}&=\frac{\sum_t\gamma_t(j,m)\bm{x}_t}{\sum_t\gamma_t(j,m)},\qquad
 c_{j,m}=\frac{\sum_t\gamma_t(j,m)}{\sum_t\gamma_t(j)},\\
 \bm{\sigma}_{j,m}^2&=\frac{\sum_t\gamma_t(j,m)(\bm{x}_t-\bm{\mu}_{j,m})^{\odot2}}{\sum_t\gamma_t(j,m)}.
\end{aligned}
 \label{eq:nlp-speech-gmm-update}
\end{equation}
转移概率按第\ref{sec:nlp-classical-baum-welch}小节的期望转移计数更新。实际流程从单音素模型开始，先把各段均匀切分作为初始对齐，迭代收敛后再展开为三音素、建决策树并继续迭代。

HMM--GMM的优点是每一帧属于哪个状态都可以检查，对齐和时长可解释，所需数据量也较小。它的局限来自两个假设：给定状态后各帧独立，相邻帧的强相关被忽略；GMM在对角协方差下难以刻画复杂的特征分布，也难以直接使用多帧拼接的高维输入。

\subsubsection{混合HMM--DNN与解码图}
\label{sec:nlp-speech-hybrid}

GMM表达力的不足可以用判别式神经网络弥补，同时保留HMM的状态结构和解码流程。混合系统的深度网络读取以第$t$帧为中心的一个窗口，例如前后各5帧、每帧40维log-Mel，拼成440维输入，经过多个隐藏层后对数千个绑定状态做softmax，输出状态后验$p(s\mid\bm{x}_t)$\scite{nlp-speech-hinton2012}。帧级训练标签来自已训练好的HMM--GMM在转写上的强制对齐（第\ref{sec:nlp-speech-forced-alignment}小节），损失为逐帧交叉熵；之后还可以用整句的序列判别目标继续训练。

HMM的解码需要似然$p(\bm{x}_t\mid s)$，网络给出的却是后验。由贝叶斯公式，
\begin{equation}
 \log p(\bm{x}_t\mid s)=\log p(s\mid\bm{x}_t)-\log p(s)+\log p(\bm{x}_t).
 \label{eq:nlp-speech-hybrid-likelihood}
\end{equation}
最后一项对同一帧的所有状态相同，比较候选路径时可以略去；状态先验$p(s)$用对齐中各状态出现的频率估计。因此，混合系统用后验除以先验得到的伪似然替换式\eqref{eq:nlp-speech-gmm}，其余部分不变。网络后验与HMM转移、语言模型分数的尺度不一致，解码时还要给声学分数乘一个小于1的缩放系数，在开发集上选定。

接下来要把HMM状态、三音素、词典和语言模型组合成一张可以搜索的图。\term{加权有限状态转导器}（weighted finite-state transducer, WFST）提供了统一的表示：它是一个有向图，每条弧带有输入符号、输出符号和权重，一条从初态到终态的路径把它读入的输入串映射成输出串，路径权重为各弧权重之和（权重取负对数概率）。两个转导器$T_1$与$T_2$的\term{复合}（composition）$T_1\circ T_2$把$T_1$的输出接到$T_2$的输入：若$T_1$把$a$映射为$b$、代价为$w_1$，$T_2$把$b$映射为$c$、代价为$w_2$，则复合结果把$a$直接映射为$c$，代价为对所有中间串$b$取$\min(w_1+w_2)$。复合的含义就是把两步映射合成一步，并把两步的代价相加。

识别中依次使用四个转导器：$H$把HMM转移标识映射为上下文相关音素，其中编码了各状态的拓扑与转移代价；$C$把上下文相关音素映射为上下文无关音素；$L$是发音词典，把音素串映射为词；$G$是语言模型，作为只接受词序列的接受器，权重为$n$元模型的负对数概率。Kaldi工具包把它们组合为
\begin{equation}
 \mathit{HCLG}=\min\bigl(\det(H\circ C\circ L\circ G)\bigr),
 \label{eq:nlp-speech-hclg}
\end{equation}
这里略去了实现中逐级插入的确定化与最小化步骤\scite{nlp-speech-mohri2002}。\term{确定化}（determinization）使每个状态出发的弧不出现相同的输入符号，每个输入串至多对应一条路径，搜索时不必在等价路径上重复计算；\term{最小化}（minimization）合并后续行为完全相同的状态，缩小图的规模。词典中的同音词会使$L$无法确定化：若“hi”和“high”都读作/h ay/，同一音素串对应两个词。实现时在同音词的音素串末尾加上不同的消歧符号，确定化后再去掉。复合后的图输入转移标识、输出词，权重中已经包含转移与语言模型分数。

解码在$\mathit{HCLG}$上做带剪枝的Viterbi搜索。每读入一帧，活动路径沿输入非空的弧前进一步，代价加上$-\kappa\log p(\bm{x}_t\mid s)$，其中$s$由弧上的转移标识确定，$\kappa$是上述声学缩放系数；代价比当前最优路径高出束宽以上的路径被删除。经过输出词的弧时记录该词，最后回溯得到词序列，也可以保留一个词格供第\ref{sec:nlp-speech-fusion}小节重打分\scite{nlp-speech-kaldi}

这条路线的错误可以分成三类定位。发音词典没有“海棠楼”时，解码图里根本没有这个词，无论声学模型多好都无法输出，属于词典缺项；声学模型给错误状态更高的分数，属于评分错误；正确路径的总分其实更高，却在搜索中被剪掉，属于搜索错误。区分后两类的办法是计算参考文本在图中的强制路径分数：若它比输出路径分数高，就是搜索错误，应放宽束宽而不是重训模型。混合系统的主要负担在工程上：单音素、三音素、决策树、特征变换、说话者自适应和神经网络各阶段都需要单独训练和调参，前一阶段的对齐错误还会传给后一阶段。

\subsection{对齐路径求和：CTC}
\label{sec:nlp-speech-ctc}

混合系统仍依赖帧级对齐标签和多个独立训练的阶段。如果只有声音与整句文字，又不想先指定一个唯一的帧级对齐，可以把所有合法对应的概率加起来作为文本的概率。\term{连接时序分类}（connectionist temporal classification, CTC）就是这样定义似然的\scite{nlp-speech-graves2006}。它要求每个编码步的输出在给定声音后相互条件独立，因此结构简单、天然单调、容易流式；代价是不显式使用输出历史，常需外部语言模型补足拼写与语言规律。CTC在全书其他位置没有正文，本小节完整写出递推、梯度与解码。

\subsubsection{扩展标签与前向--后向递推}
\label{sec:nlp-speech-ctc-forward-backward}

CTC在词表$V$外增加\term{空白符}（blank）$\varnothing$，编码器之后接一个线性层与softmax，在每个编码步给出$|V|+1$维分布$q_t(k)=p(k\mid\bm{X},t)$。空白表示此步不输出新符号，它不等于静音。长度为$T$的路径$\pi=(\pi_1,\ldots,\pi_T)$每步取一个符号；记$\mathcal{B}$为“先合并相邻相同符号，再删除空白”的折叠变换，则
\begin{equation}
 p(y\mid\bm{X})=
 \sum_{\pi:\mathcal{B}(\pi)=y}\prod_{t=1}^{T}q_t(\pi_t),
 \qquad L_{\mathrm{CTC}}=-\log p(y\mid\bm{X}).
 \label{eq:nlp-speech-ctc-probability}
\end{equation}
乘积形式假定给定$\bm{X}$后各步的符号选择条件独立。这里独立的是路径上的输出，不是声音帧：双向编码器或自注意力编码器仍可以让每个$q_t$依赖整段声音。但$q_t$不读取已经输出的文字，模型无法直接学到“上一个字是‘四’，下一个字更可能是‘月’”这类规律。

折叠的顺序不能颠倒：路径$(a,a,\varnothing,a)$先合并为$(a,\varnothing,a)$，再删空白得到$aa$；若先删空白再合并，就会错误地得到$a$。因此目标中相邻的相同符号之间至少要有一个空白隔开。令$r(y)=\sum_{u=2}^{U}\mathbf{1}[y_u=y_{u-1}]$为相邻重复数，非空目标存在合法路径的条件是$T\geq U+r(y)$。这是对编码器输出长度的约束：第\ref{sec:nlp-speech-encoder}小节的降采样倍数过大时，这个条件可能不满足，样本的损失为无穷大。

直接枚举路径的数量随$T$指数增长，需要动态规划。在目标首尾和相邻符号之间插入空白，得到\term{扩展标签}（extended label）
$z=(\varnothing,y_1,\varnothing,\ldots,y_U,\varnothing)$，
长度为$2U+1$。定义前向变量$\alpha_{t,s}$为前$t$步、终于扩展标签位置$s$的所有合法部分路径的概率之和。对$U>0$，初值为
\begin{equation}
 \alpha_{1,1}=q_1(\varnothing),\quad
 \alpha_{1,2}=q_1(y_1),\quad
 \alpha_{1,s}=0\quad(s>2).
 \label{eq:nlp-speech-ctc-initial}
\end{equation}
越界位置的概率取0。到达位置$s$有三种来源：停在原位置，从前一位置进入，或跳过中间的空白从$s-2$进入。最后一种只在当前标签非空白、且与$s-2$处的标签不同时允许，否则目标$aa$中的两个$a$会被走成一个持续的$a$：
\begin{equation}
\begin{aligned}
 \delta_s&=\mathbf{1}[s>2,\ z_s\ne\varnothing,\ z_s\ne z_{s-2}],\\
 \alpha_{t,s}
 &=q_t(z_s)\bigl(\alpha_{t-1,s}+\alpha_{t-1,s-1}
                    +\delta_s\alpha_{t-1,s-2}\bigr),\\
 p(y\mid\bm{X})&=\alpha_{T,2U}+\alpha_{T,2U+1}.
\end{aligned}
\label{eq:nlp-speech-ctc-forward}
\end{equation}
合法的结束位置只有最后一个标签和尾部空白两个。空目标只有全空白路径，概率为$\prod_tq_t(\varnothing)$。

表\ref{tab:nlp-speech-ctc-numbers}给出三步完整分布及目标$ab$的前向变量，分布为教学构造。第三步到达位置4时，三个来源的前向变量为$0.04,0.16,0.45$，乘以$q_3(b)=0.1$得$0.065$；再加尾空白的$0.020$，目标概率为$0.085$。

\begin{table}[htbp]
  \centering
  \begin{tabular}{@{}crrrrrrrr@{}}
    \toprule
    & \multicolumn{3}{c}{逐步分布} & \multicolumn{5}{c}{$z=(\varnothing,a,\varnothing,b,\varnothing)$}\\
    $t$ & $\varnothing$ & $a$ & $b$ & $\alpha_{t,1}$ & $\alpha_{t,2}$ & $\alpha_{t,3}$ & $\alpha_{t,4}$ & $\alpha_{t,5}$\\
    \midrule
    1 & .5 & .4 & .1 & .5 & .4 & 0 & 0 & 0\\
    2 & .4 & .5 & .1 & .2 & .45 & .16 & .04 & 0\\
    3 & .5 & .4 & .1 & .1 & .26 & .305 & .065 & .020\\
    \bottomrule
  \end{tabular}
  \caption{CTC前向变量的教学计算。逐步分布为构造数值，前向变量由式\eqref{eq:nlp-speech-ctc-forward}递推得到；每行的前向变量是部分路径概率之和，不是归一化分布。}
  \label{tab:nlp-speech-ctc-numbers}
\end{table}

同一结果可以完全枚举核对。折叠为$ab$的五条路径是
$(a,a,b)$、$(a,b,b)$、$(a,\varnothing,b)$、$(\varnothing,a,b)$和$(a,b,\varnothing)$，
概率分别为$0.020,0.004,0.016,0.025,0.020$，总和$0.085$，损失$-\log0.085=2.4651$。目标$aa$在这三步中只有$(a,\varnothing,a)$一条路径，概率为$0.064$。

求梯度还需要后向变量。定义$\beta_{t,s}$为第$t$步位于位置$s$时，后续第$t+1$至$T$步能够合法走到终点的所有路径后缀的概率之和；它不含第$t$步本身的输出概率。终点条件和逆序递推为
\begin{equation}
\begin{aligned}
 \beta_{T,s}&=\mathbf{1}[s\in\{2U,2U+1\}],\\
 \beta_{t,s}&=q_{t+1}(z_s)\beta_{t+1,s}+q_{t+1}(z_{s+1})\beta_{t+1,s+1}
   +\delta_{s+2}\,q_{t+1}(z_{s+2})\beta_{t+1,s+2},
\end{aligned}
\label{eq:nlp-speech-ctc-backward}
\end{equation}
三项与前向递推的三种来源一一对应，越界项取0。按这个定义，$\alpha_{t,s}\beta_{t,s}$恰好是所有在第$t$步经过位置$s$的完整路径概率之和。每条合法路径在每一步都恰好经过一个位置，于是对任意$t$都有$\sum_s\alpha_{t,s}\beta_{t,s}=p(y\mid\bm{X})$，这是检查实现的常用等式。在上例中，$\beta_{3,:}=(0,0,0,1,1)$，$\beta_{2,:}=(0,.1,.1,.6,.5)$，$\beta_{1,:}=(.05,.15,.10,.26,.20)$。第2步的检查为$.45\times.1+.16\times.1+.04\times.6=0.085$，第1步为$.5\times.05+.4\times.15=0.085$。

实际计算在对数域进行，以免长序列的连乘下溢。令$\ell_{t,s}=\log\alpha_{t,s}$，把式\eqref{eq:nlp-speech-ctc-forward}中的加法换成
$\operatorname{LSE}(v_1,\ldots,v_k)=m+\log\sum_j\exp(v_j-m)$，
其中$m=\max_jv_j$；不可达项取$-\infty$，全部不可达时直接返回$-\infty$。后向变量同样处理。前向与后向各需$O(TU)$次状态更新；只求损失时可以滚动保存两行，求梯度则要保存或重算全部状态。

\subsubsection{损失梯度与尖峰现象}
\label{sec:nlp-speech-ctc-gradient}

有了前向与后向变量，就能算出每一步每个符号被“占用”的后验概率。第$t$步位于位置$s$的后验为$\alpha_{t,s}\beta_{t,s}/p(y\mid\bm{X})$；把同一符号的所有位置加起来，得到符号$k$在第$t$步的\term{后验占用}（posterior occupancy）
\begin{equation}
 \gamma_t(k)=\frac{1}{p(y\mid\bm{X})}\sum_{s:\,z_s=k}\alpha_{t,s}\beta_{t,s},
 \qquad \sum_{k}\gamma_t(k)=1.
 \label{eq:nlp-speech-ctc-occupancy}
\end{equation}
它表示：在所有能折叠为目标的路径中按概率加权，第$t$步输出$k$的比例。

设第$t$步softmax之前的未归一化分数为$o_{t,k}$，$q_t=\operatorname{softmax}(\bm{o}_t)$。$q_t(k)$在每条经过$(t,s)$且$z_s=k$的路径中恰好出现一次，因此$\partial p/\partial q_t(k)=\sum_{s:z_s=k}\alpha_{t,s}\beta_{t,s}/q_t(k)$，即$\partial\log p/\partial q_t(k)=\gamma_t(k)/q_t(k)$。再用softmax的导数$\partial q_t(j)/\partial o_{t,k}=q_t(j)(\mathbf{1}[j=k]-q_t(k))$求和：
\begin{equation}
 \frac{\partial L_{\mathrm{CTC}}}{\partial o_{t,k}}
 =-\sum_j\frac{\gamma_t(j)}{q_t(j)}\,q_t(j)\bigl(\mathbf{1}[j=k]-q_t(k)\bigr)
 =q_t(k)-\gamma_t(k).
 \label{eq:nlp-speech-ctc-gradient}
\end{equation}
最后一步用到$\sum_j\gamma_t(j)=1$。这个结果与普通分类的交叉熵梯度形式相同，只是“标签”换成了由当前模型自己算出的后验占用：CTC相当于在每一步用软目标$\gamma_t$训练分类器，而软目标随模型更新而改变。这与EM的思路一致，E步算对齐的后验，M步向它拟合，区别在于CTC用梯度一步完成两者。

在表\ref{tab:nlp-speech-ctc-numbers}的例子中，第2步的位置后验为$(0,.529,.188,.282,0)$，于是$\gamma_2(\varnothing)=.188$、$\gamma_2(a)=.529$、$\gamma_2(b)=.282$。梯度$q_2-\gamma_2=(.212,-.029,-.182)$：梯度下降会降低这一步空白的分数，主要提高$b$的分数。

这个梯度也解释了训练后CTC输出的\term{尖峰现象}（peaky behavior）：多数步以接近1的概率输出空白，每个非空白符号只在一两步上出现高概率的尖峰\scite{nlp-speech-zeyer2021}。原因在于两类符号在路径中的地位不同。每个标签至少占一步，而空白可以填充标签之间任意多步，并且扩展标签中有$U+1$个空白位置。训练初期各符号概率接近均匀，后验占用中空白占很大比例，梯度便在多数步上推高空白。空白概率升高后，每个标签只需在其相对概率最高的一步出现即可完成折叠，占用进一步集中到这一步，形成正反馈。尖峰只表示“在这一步输出该符号”，CTC目标从未要求它落在发音的起点、中心或终点；单向流式编码器还常把尖峰推迟到看见足够声音之后。图\ref{fig:nlp-speech-alignment}用两条折叠为同一文本的路径说明，符号出现在哪一步本身就是不确定的。

\begin{figure*}[htbp]
  \centering
  % Adapted from the academic-drawing TikZ template; dimensions are final size.
\begin{tikzpicture}[
  x=1mm,y=1mm,
  every node/.style={font=\figurefont\small,align=center,text=BookInk},
  cell/.style={draw=BookRule,line width=.8pt,minimum width=13mm,
    minimum height=9mm,inner sep=1mm},
  label/.style={anchor=east},
  flow/.style={book arrow}
]
  \foreach \i in {1,...,6} {
    \node[cell,fill=FigureInput!12] at ({27+17*\i},0) {$\bm{h}_{\i}$};
  }
  \node[label] at (31,0) {编码步};
  \node[label] at (31,-15) {路径A};
  \node[label] at (31,-30) {路径B};
  \foreach \i/\v in {1/八,2/八,3/{$\varnothing$},4/八,5/{$\varnothing$},6/{$\varnothing$}} {
    \node[cell,fill=FigureModel!10] at ({27+17*\i},-15) {\v};
  }
  \foreach \i/\v in {1/{$\varnothing$},2/八,3/{$\varnothing$},4/八,5/八,6/{$\varnothing$}} {
    \node[cell,fill=FigureModel!10] at ({27+17*\i},-30) {\v};
  }
  \draw[flow] (137,-15) -- (145,-15);
  \draw[flow] (137,-30) -- (145,-30);
  \node[anchor=west,text=FigureOutput] at (146,-15) {八八};
  \node[anchor=west,text=FigureOutput] at (146,-30) {八八};
  \draw[BookRule,line width=.8pt] (5,-42) -- (160,-42);
  \node[label] at (31,-54) {独立标注};
  \node[anchor=west,align=left] at (37,-54)
    {第一个“八”：0.02--0.10\,s，说话者甲\\
     第二个“八”：0.13--0.22\,s，说话者甲};
\end{tikzpicture}

  \caption{同一声学序列的两条CTC教学路径均折叠为“八八”，但符号出现的位置不同。空白为$\varnothing$，相邻重复先合并，再删除空白。最下方的时间区间与说话者是另行给定的示意标注，不是由两条路径证明的精确时间。}
  \label{fig:nlp-speech-alignment}
\end{figure*}
\FloatBarrier

\subsubsection{贪心解码与前缀束搜索}
\label{sec:nlp-speech-ctc-decoding}

最简单的解码是每步取概率最大的符号再折叠。但这样得到的是概率最大的单条路径所对应的文本，而不是总概率最大的文本，因为一个文本的概率是许多路径之和。构造两步相同分布$q(\varnothing)=0.40$、$q(a)=0.35$、$q(b)=0.25$。贪心选出全空白路径，对应空文本，概率$0.16$；文本$a$却汇总了$(a,a)$、$(a,\varnothing)$、$(\varnothing,a)$三条路径，概率$0.1225+0.14+0.14=0.4025$，远高于空文本。

前缀束搜索把折叠后对应同一文字前缀的路径合并，但必须区分路径的末步是空白还是非空白。令$b_t(l)$与$n_t(l)$分别为前$t$步折叠为$l$、且末步为空白与非空白的概率之和；初值$b_0(\epsilon)=1$，其余为0。区分两者的原因是：若末字为$a$的前缀再遇到$a$，非空白结尾的路径只能延续同一个$a$，空白结尾的路径才能形成新的$a$。算法\ref{alg:nlp-speech-prefix}给出这一过程\scite{nlp-speech-hannun2014}

\begin{algorithm}[CTC前缀束搜索]
  \label{alg:nlp-speech-prefix}
  \begin{algorithmic}[1]
    \Require 分布$q_1,\ldots,q_T$，词表$V$，束宽$K$
    \Ensure 按合并概率选择的文字序列
    \State 初始化$b_0(\epsilon)=1,\ n_0(\epsilon)=0$，束为$\{\epsilon\}$
    \For{$t=1,\ldots,T$}
      \State 新映射$b',n'$全部置0
      \For{旧束中的每个前缀$l$}
        \State $s\gets b_{t-1}(l)+n_{t-1}(l)$
        \State $b'(l)\mathrel{+}=s\,q_t(\varnothing)$ \Comment{补一个空白，前缀不变}
        \For{$c\in V$}
          \If{$l\ne\epsilon$且$c=\operatorname{last}(l)$}
            \State $n'(l)\mathrel{+}=n_{t-1}(l)q_t(c)$ \Comment{延续同一个字}
            \State $n'(lc)\mathrel{+}=b_{t-1}(l)q_t(c)$ \Comment{空白之后形成新字}
          \Else
            \State $n'(lc)\mathrel{+}=s\,q_t(c)$
          \EndIf
        \EndFor
      \EndFor
      \State 合并同名前缀后，按$b'(l)+n'(l)$保留前$K$项，赋给$b_t,n_t$
    \EndFor
    \State \Return $\argmax_l\{b_T(l)+n_T(l)\}$
  \end{algorithmic}
\end{algorithm}

对上述两步例子，前缀$a$在第一步为$(b_1,n_1)=(0,0.35)$；第二步$b_2(a)=0.35\times0.40=0.14$，$n_2(a)=0.35\times0.35+0.40\times0.35=0.2625$，合计$0.4025$，与枚举一致。不剪枝时算法精确地合并了所有前缀，保留前$K$项后才成为近似搜索。

外部语言模型在形成新前缀的位置加入。把算法中更新$n'(lc)$的两行里的$q_t(c)$换成$q_t(c)\,p_{\mathrm{lm}}(c\mid l)^{\lambda}e^{\gamma}$，其中$\lambda$是语言模型权重，$\gamma$是补偿语言模型偏好短文本的插入奖励；延续同一个字的那一行不乘语言模型因子。同一个新前缀$lc$可能在不同的步上形成，但每次乘的都是同一个因子，所以等价于对整句文本只计一次语言模型分数。第\ref{sec:nlp-speech-fusion}小节在更一般的形式下讨论这种融合。

CTC用路径求和去掉了帧级对齐标签和多阶段训练，训练只需整句转写，输出天然单调，解码可以逐步进行。它的弱点同样来自条件独立假设：每一步的输出不读取已输出文字，同音字、拼写一致性等依赖输出历史的规律只能靠编码器从声音中间接学到，或交给外部语言模型。

\subsection{注意力编码器—解码器}
\label{sec:nlp-speech-aed}

要让后一个字的概率直接依赖前面已经输出的字，可以把识别看作以声音为条件的序列生成。\term{注意力编码器—解码器}（attention-based encoder--decoder, AED）由编码器读入整段声音，解码器在每一步读取已输出的前缀，并用注意力从编码序列中取出与当前输出相关的上下文，直接最大化整句的条件概率。这种结构与第\ref{sec:nlp-generation-encoder-decoder}小节的机器翻译模型相同，教师强制训练与束搜索的一般机制分别见第\ref{sec:nlp-generation-training}小节与第\ref{sec:nlp-generation-beam}小节。它解决了CTC的条件独立问题，但带来两项新代价：注意力没有单调约束，长语音上可能跳读、重复或过早结束；注意力需要整段编码，难以流式运行。

\subsubsection{LAS金字塔编码与注意力解码}
\label{sec:nlp-speech-las}

Listen, Attend and Spell（LAS）是这类模型的代表，由编码器listener和解码器speller组成\scite{nlp-speech-chan2015}。语音特征序列比文字长得多：10\,ms帧移下，10秒声音有1000帧。注意力每步要在全部编码位置上计算权重，位置太多时不仅计算量大，训练也难以收敛。listener因此采用\term{金字塔双向LSTM}（pyramidal bidirectional LSTM, pBLSTM）：每一层把下层相邻两步的输出拼接后作为本层的一步输入，使长度减半。记第$j$层第$i$步的输出为$\bm{h}^{(j)}_i$，则
\begin{equation}
 \bm{h}^{(j)}_i=\operatorname{BLSTM}\bigl(\bm{h}^{(j)}_{i-1},\,
   [\bm{h}^{(j-1)}_{2i-1};\bm{h}^{(j-1)}_{2i}]\bigr).
 \label{eq:nlp-speech-pblstm}
\end{equation}
LSTM与双向循环网络的结构可参考\BookChapterRef{chap:recurrent-networks}{卷五循环神经网络章}。设输入为$\bm{X}\in\mathbb{R}^{T_0\times80}$，底层BLSTM每个方向$d$维，输出$T_0\times2d$；第一层金字塔把相邻两步拼成$4d$维输入，输出$(T_0/2)\times2d$。叠加三层后长度为$T_0/8$，每步对应80\,ms；10秒声音的1000帧缩为125步。降低长度同时让每个编码步汇总更长一段声音，注意力更容易找到对应位置。

speller是一个LSTM解码器。第$u$步的状态$\bm{s}_u$由前一状态、前一个输出字符$y_{u-1}$和前一步上下文$\bm{c}_{u-1}$更新；然后用$\bm{s}_u$与每个编码步打分，得到注意力权重与上下文向量：
\begin{equation}
 e_{u,t}=\langle\phi(\bm{s}_u),\psi(\bm{h}_t)\rangle,\qquad
 a_{u,t}=\frac{\exp e_{u,t}}{\sum_{t'}\exp e_{u,t'}},\qquad
 \bm{c}_u=\sum_{t}a_{u,t}\bm{h}_t,
 \label{eq:nlp-speech-las-attention}
\end{equation}
其中$\phi,\psi$是两个多层感知机，把解码状态与编码表示映射到同一空间后做内积。输出层以$[\bm{s}_u;\bm{c}_u]$经多层感知机与softmax给出字符分布。目标序列末尾加结束符，损失为
\begin{equation}
 L_{\mathrm{AED}}=-\sum_{u=1}^{U+1}
       \log p_\theta(y_u\mid\bm{X},y_{<u}),
 \qquad y_{U+1}=\texttt{EOS}.
 \label{eq:nlp-speech-las}
\end{equation}
它与式\eqref{eq:nlp-speech-ctc-probability}的根本区别在于每一项都条件于$y_{<u}$，没有任何对齐变量需要求和：对齐被注意力权重隐式地处理了。

内积打分只看内容，不知道上一步读到了哪里。语音的对应是单调的，相邻输出读取的位置应该相邻，但式\eqref{eq:nlp-speech-las-attention}并不保证这一点。\term{位置感知注意力}（location-aware attention）把上一步的注意力权重$\bm{a}_{u-1}$经一维卷积得到每个位置的特征$\bm{f}_{u,t}$，加入加性打分$e_{u,t}=\bm{w}^{\top}\tanh(\bm{W}\bm{s}_{u-1}+\bm{V}\bm{h}_t+\bm{U}\bm{f}_{u,t}+\bm{b})$\scite{nlp-speech-chorowski2015}。加性打分的形式与式\eqref{eq:nlp-generation-bahdanau}相同，新增的$\bm{f}_{u,t}$让模型知道上一步读到了哪里，倾向于在附近继续读取。这一做法减轻了跳读，但它仍是学到的倾向，而不是硬约束。

仍用“八八”短录音，两个“八”由两个输出步分别生成，不需要CTC的空白分隔。假设候选“八／八／EOS”的三个条件概率为$0.7,0.6,0.8$，候选概率为$0.336$；另一候选必须按自己的前缀重新计算各步条件概率，不能只替换一个字的概率。训练时教师强制让每一步读取正确前缀，推理时却只能读取自己生成的前缀，这种不一致使早期错误可能连锁传播；LAS原文在训练中以一定概率用模型自己的采样代替正确前缀来缓解这一点。

AED的两项代价都来自注意力的全局性。注意力权重可以落在任何位置，模型在长语音上可能跳过一段、反复读同一段，或者在说完之前就输出EOS；训练集中没有的长度尤其容易出问题。双向编码与全局注意力又要求整段声音都已到达，不能边听边出字。

\subsubsection{CTC与注意力的联合训练}
\label{sec:nlp-speech-joint-ctc-attention}

CTC有单调约束但条件独立，AED读取历史但没有单调约束，二者的优缺点恰好互补。联合训练在共享编码器上同时接两个输出分支：一个线性层加softmax的CTC分支，以及一个注意力解码器，损失取加权和\scite{nlp-speech-kim2017}：
\begin{equation}
 L=\lambda L_{\mathrm{CTC}}+(1-\lambda)L_{\mathrm{AED}},\qquad 0\leq\lambda\leq1.
 \label{eq:nlp-speech-joint-loss}
\end{equation}
$\lambda$在开发集上选取。CTC损失只在单调路径上有概率，它迫使共享编码器的表示与文字顺序对齐，注意力解码器因而更快学到合理的对齐，训练早期收敛也更快。

解码时也可以让两个分支共同打分\scite{nlp-speech-watanabe2017}。注意力解码器按束搜索逐字扩展前缀，每个前缀同时计算CTC分支给出的\term{CTC前缀分数}（CTC prefix score），即所有折叠后以该前缀开头的文本的总概率。对前缀$h=gc$，$c$是新加入的字，按$c$第一次出现的步$t$分类求和：
\begin{equation}
 \Psi(gc)=\sum_{t=1}^{T}\Phi_{t-1}(g)\,q_t(c),\qquad
 \Phi_{t-1}(g)=\begin{cases}
   b_{t-1}(g), & c=\operatorname{last}(g),\\
   b_{t-1}(g)+n_{t-1}(g), & \text{其他},
 \end{cases}
 \label{eq:nlp-speech-ctc-prefix-score}
\end{equation}
其中$b,n$是算法\ref{alg:nlp-speech-prefix}中的空白结尾与非空白结尾概率，按完整的$T$步计算而不剪枝。若候选在$g$之后输出EOS，CTC分数取整句概率$b_T(g)+n_T(g)$。联合分数为
\begin{equation}
 S(h)=\lambda\log\Psi(h)+(1-\lambda)\log p_{\mathrm{AED}}(h\mid\bm{X})
      +\mu\log p_{\mathrm{lm}}(h),
 \label{eq:nlp-speech-joint-score}
\end{equation}
最后一项是可选的外部语言模型。仍用两步分布$q(\varnothing)=.40$、$q(a)=.35$、$q(b)=.25$：$\Psi(a)=1\times.35+.40\times.35=.49$，它等于以$a$开头的两个文本$a$（$.4025$）与$ab$（$.0875$）的概率之和；若候选在$a$之后结束，CTC分数取$.4025$。

CTC前缀分数的作用在于否决不单调的候选。若注意力解码器跳过了一段声音，跳过的字无法与CTC的单调路径对上，$\Psi$会很小；若解码器过早输出EOS，剩余声音找不到对应的文字，整句CTC概率也很小。另一种做法是先用其中一个分支生成$n$个完整候选，再用另一个分支重打分，实现简单，但只能在已有候选中选择。联合模型的代价是两个分支都要计算，且解码仍需完整编码，并未解决流式问题；第\ref{sec:nlp-speech-chunk-attention}小节的两遍解码把它与流式编码器结合起来。

\subsection{历史条件与路径求和：RNN-T}
\label{sec:nlp-speech-rnnt}

CTC可以流式但不读输出历史，AED读取历史却不单调、难以流式。\term{循环神经网络转导器}（recurrent neural network transducer, RNN-T）把两者的优点放进同一个概率模型\scite{nlp-speech-graves2012}。它在声学步$t$与输出位置$u$组成的二维格点上定义两类边：空白边推进时间，标签边推进输出，同一个声学步可以连续输出多个符号。每个格点的输出分布同时依赖编码器的$\bm{h}_t$和已输出的前缀，路径求和仍保证单调对齐。代价是联合网络要在$T\times U$个格点上各计算一次词表分布，训练显存很大。

\subsubsection{预测网络与联合网络}
\label{sec:nlp-speech-rnnt-networks}

RNN-T由三部分组成，结构见图\ref{fig:nlp-speech-rnnt-structure}。声学编码器与CTC相同，把$\bm{X}$变成$\bm{h}_1,\ldots,\bm{h}_T$，默认采用第\ref{sec:nlp-speech-encoder}小节的Conformer。\term{预测网络}（prediction network）只读取已输出的非空白符号：以开始符号$y_0$和$y_1,\ldots,y_u$为输入，经嵌入层和一两层LSTM得到$\bm{g}_u$。它的作用类似一个只看文字的语言模型，为下一个符号提供历史条件。\term{联合网络}（joint network）把两路表示合并：
\begin{equation}
 \bm{z}_{t,u}=\bm{W}_o\tanh(\bm{W}_a\bm{h}_t+\bm{W}_p\bm{g}_u+\bm{b}),\qquad
 q_{t,u}=\operatorname{softmax}(\bm{z}_{t,u})\in\mathbb{R}^{|V|+1},
 \label{eq:nlp-speech-rnnt-joint}
\end{equation}
其中$\bm{W}_a\in\mathbb{R}^{d_{\mathrm{joint}}\times d_{\mathrm{enc}}}$、$\bm{W}_p\in\mathbb{R}^{d_{\mathrm{joint}}\times d_{\mathrm{pred}}}$把两路投影到同一维度，$\bm{W}_o\in\mathbb{R}^{(|V|+1)\times d_{\mathrm{joint}}}$输出含空白的词表分数。Graves的原始形式直接把两路输出的分数相加再做softmax，式\eqref{eq:nlp-speech-rnnt-joint}在相加之后多了一层非线性。

\begin{figure}[htbp]
  \centering
  % 第九卷 15.2.5.1 节：RNN-T的编码器、预测网络与联合网络（正文栏宽，约108 mm）。
% 蓝色：输入；黄色描边：网络模块；红色：输出分布。虚线为推断时的反馈。
\begin{tikzpicture}[
  x=1mm,y=1mm,
  figure picture,
  rinp/.style={figure node,draw=FigureInput,fill=FigureInputFill,
    minimum width=40mm,minimum height=8mm,inner ysep=2pt,font=\figurefont\scriptsize},
  rmod/.style={figure node,draw=FigureModel,fill=FigureModelFill,
    minimum width=40mm,minimum height=10mm,inner ysep=2pt,font=\figurefont\scriptsize},
  rout/.style={figure node,draw=FigureOutput,fill=FigureOutputFill,
    minimum width=86mm,minimum height=8mm,inner ysep=2pt,font=\figurefont\scriptsize},
  fl/.style={figure flow},
  dl/.style={figure flow,draw=FigureMuted,dash pattern=on 2pt off 1.5pt},
  el/.style={font=\figurefont\scriptsize,text=FigureMuted,inner sep=1pt,align=center}
]
  \node[rinp] (x) at (22,0) {声学特征$\bm{X}$};
  \node[rinp] (y) at (72,0) {已输出的非空白符号$y_0,y_1,\ldots,y_u$};
  \node[rmod] (enc) at (22,15) {声学编码器（如Conformer）\\$\bm{H}\in\mathbb{R}^{B\times T\times d_{\mathrm{enc}}}$};
  \node[rmod] (pred) at (72,15) {预测网络（嵌入＋LSTM）\\$\bm{G}\in\mathbb{R}^{B\times(U+1)\times d_{\mathrm{pred}}}$};
  \node[rmod,minimum width=86mm] (joint) at (47,33)
    {联合网络：$\bm{z}_{t,u}=\bm{W}_o\tanh(\bm{W}_a\bm{h}_t+\bm{W}_p\bm{g}_u)$\\
     按$(t,u)$广播相加，隐藏张量$B\times T\times(U+1)\times d_{\mathrm{joint}}$};
  \node[rout] (q) at (47,49) {softmax：$q_{t,u}=\operatorname{softmax}(\bm{z}_{t,u})$，
    张量$B\times T\times(U+1)\times(|V|+1)$};

  \draw[fl] (x) -- (enc);
  \draw[fl] (y) -- (pred);
  \draw[fl] (enc) -- node[el,left] {$\bm{h}_t$} (enc |- joint.south);
  \draw[fl] (pred) -- node[el,right] {$\bm{g}_u$} (pred |- joint.south);
  \draw[fl] (joint) -- (q);
  \draw[dl] (q.east) -- ++(6,0) |- node[el,pos=.25,right] {推断时\\非空白\\输出送回} (y.east);
\end{tikzpicture}

  \caption{RNN-T的三部分结构。编码器只读声音，预测网络只读已输出的非空白符号，联合网络把$\bm{h}_t$与$\bm{g}_u$投影后相加、经$\tanh$与输出层得到格点$(t,u)$上含空白的分布。训练时$\bm{H}$与$\bm{G}$按$(t,u)$广播成四维张量；推断时非空白输出沿虚线送回预测网络，空白则让编码器前进一步。}
  \label{fig:nlp-speech-rnnt-structure}
\end{figure}

张量形状说明了计算量落在哪里。批大小为$B$时，编码器输出$\bm{H}\in\mathbb{R}^{B\times T\times d_{\mathrm{enc}}}$，预测网络输出$\bm{G}\in\mathbb{R}^{B\times(U+1)\times d_{\mathrm{pred}}}$，多出的一行对应只输出了开始符号的位置$u=0$。两路投影后分别扩展为$B\times T\times1\times d_{\mathrm{joint}}$与$B\times1\times(U+1)\times d_{\mathrm{joint}}$，广播相加得到$B\times T\times(U+1)\times d_{\mathrm{joint}}$，最后输出$B\times T\times(U+1)\times(|V|+1)$。编码器和预测网络都只算一次，四维张量完全来自联合网络。

预测网络不读声音，是有意的设计。这样声学信息只经编码器进入，编码器可以独立地做成流式；预测网络也可以像语言模型一样在推断时缓存状态，只在输出非空白符号时更新。由于空白只推进时间、不进入预测网络，RNN-T的路径恢复只需删除空白，不需要合并重复：目标$aa$可以在同一时间步连续发出两个$a$，CTC的长度条件$T\geq U+r(y)$在这里不适用。若令$\bm{W}_p=\bm{0}$，分布不再依赖输出历史，模型就退化为一种条件独立的转导器，与CTC相近，但同一步仍允许输出多个符号，两者并不等同。实践中预测网络甚至可以只看最近一两个符号，称为无状态预测网络，准确率下降不多\scite{nlp-speech-ghodsi2020}；这说明RNN-T的历史条件主要提供局部的拼写与搭配信息。

\subsubsection{格点前向--后向与损失}
\label{sec:nlp-speech-rnnt-loss}

格点上的索引取$t=1,\ldots,T$、$u=0,\ldots,U$。空白边从$(t,u)$到$(t+1,u)$，消耗一个声学步；标签$y_{u+1}$从$(t,u)$到$(t,u+1)$，增加一个输出而不推进声音。固定目标后，记$b_{t,u}=q_{t,u}(\varnothing)$、$e_{t,u}=q_{t,u}(y_{u+1})$。前向变量$\alpha_{t,u}$汇总到达格点$(t,u)$的所有部分路径概率，不含该格点本身的输出。初值$\alpha_{1,0}=1$，越界项为0：
\begin{equation}
\begin{aligned}
 \alpha_{t,u}
 &=\alpha_{t-1,u}b_{t-1,u}
   +\alpha_{t,u-1}e_{t,u-1},\\
 p(y\mid\bm{X})&=\alpha_{T,U}b_{T,U}.
\end{aligned}
\label{eq:nlp-speech-rnnt-forward}
\end{equation}
首列$u=0$只能由空白进入，首行$t=1$只能由标签进入。最后必须从$(T,U)$再发出一次空白，消耗最后一个声学步，所以终点概率要乘$b_{T,U}$。

后向变量$\beta_{t,u}$汇总从格点$(t,u)$出发、含该格点输出在内走到终点的剩余概率。初值$\beta_{T,U}=b_{T,U}$，逆序递推为
\begin{equation}
 \beta_{t,u}=b_{t,u}\beta_{t+1,u}
       +e_{t,u}\beta_{t,u+1},
 \label{eq:nlp-speech-rnnt-backward}
\end{equation}
只保留合法边：$u=U$时没有标签边，$t=T$时没有通往下一声学步的空白边。于是$\beta_{1,0}=p(y\mid\bm{X})$，而$\alpha_{t,u}\beta_{t,u}$是所有经过格点$(t,u)$的路径概率之和。

构造$T=2$、$y=aa$的例子，词表只有$a$和空白。第一列三个格点的$(b,q(a))$依次为$(.4,.6)$、$(.5,.5)$、$(.7,.3)$，第二列依次为$(.3,.7)$、$(.2,.8)$、$(.9,.1)$，合法边见图\ref{fig:nlp-speech-rnnt-lattice}。三条完整路径
$(a,a,\varnothing,\varnothing)$、
$(a,\varnothing,a,\varnothing)$、
$(\varnothing,a,a,\varnothing)$
的概率为$0.189$、$0.216$、$0.2016$，总和$0.6066$。前向递推得$\alpha_{1,:}=(1,.6,.3)$、$\alpha_{2,:}=(.4,.58,.674)$，末尾乘$.9$得到同一结果；后向递推得$\beta_{2,:}=(.504,.72,.9)$、$\beta_{1,:}=(.6066,.675,.63)$。CTC在两步内不能产生$aa$，这个例子显示了两种路径空间的区别。

\begin{figure}[htbp]
  \centering
  % Target-specific probability lattice, not a measured alignment.
\begin{tikzpicture}[
  x=1mm,y=1mm,
  every node/.style={font=\figurefont\small,text=BookInk},
  state/.style={circle,draw=FigureModel,fill=FigureModel!10,
    minimum size=5mm,inner sep=0pt,line width=.9pt},
  time/.style={book arrow,draw=BookBlue},
  token/.style={book arrow,draw=BookTeal},
  edge label/.style={fill=white,inner sep=1pt}
]
  \foreach \u in {0,1,2} {
    \node[state] (a\u) at (20,{19*\u}) {};
    \node[state] (b\u) at (62,{19*\u}) {};
    \node[anchor=east] at (9,{19*\u}) {$u=\u$};
  }
  \node at (20,-10) {$t=1$};
  \node at (62,-10) {$t=2$};
  \draw[time] (a0) -- node[edge label,above] {$\varnothing:\ .4$} (b0);
  \draw[time] (a1) -- node[edge label,above] {$\varnothing:\ .5$} (b1);
  \draw[time] (a2) -- node[edge label,above] {$\varnothing:\ .7$} (b2);
  \draw[token] (a0) -- node[edge label,right] {$a:\ .6$} (a1);
  \draw[token] (a1) -- node[edge label,right] {$a:\ .5$} (a2);
  \draw[token] (b0) -- node[edge label,right] {$a:\ .7$} (b1);
  \draw[token] (b1) -- node[edge label,right] {$a:\ .8$} (b2);
  \node[draw=FigureOutput,fill=FigureOutput!12,minimum height=8mm,
    inner sep=2mm] (end) at (97,38) {终点};
  \draw[time] (b2) -- node[edge label,above] {$\varnothing:\ .9$} (end);
  \node[anchor=west] at (5,-22) {起点$(1,0)$概率为1；终点概率为0.6066};
\end{tikzpicture}

  \caption{两编码步、目标“$aa$”的RNN-T格点。横向空白推进时间，纵向$a$推进输出；标注为教学边概率。只有从右上角发出的末尾空白进入终点，重复标签没有折叠。}
  \label{fig:nlp-speech-rnnt-lattice}
\end{figure}

梯度的推导与CTC类似。$q_{t,u}(k)$只出现在从格点$(t,u)$发出的边上，因此$\partial p/\partial q_{t,u}(k)=\alpha_{t,u}\beta'_{t,u}(k)$，其中$\beta'_{t,u}(\varnothing)=\beta_{t+1,u}$（约定$\beta_{T+1,U}=1$），$\beta'_{t,u}(y_{u+1})=\beta_{t,u+1}$，其余符号为0。经过softmax导数，得到
\begin{equation}
 \frac{\partial L}{\partial z_{t,u,k}}
 =\frac{\alpha_{t,u}\beta_{t,u}}{p}\,q_{t,u}(k)
  -\frac{\alpha_{t,u}\,q_{t,u}(k)\,\beta'_{t,u}(k)}{p}.
 \label{eq:nlp-speech-rnnt-gradient}
\end{equation}
第一项是格点占用乘以模型分布，第二项是沿符号$k$的边离开该格点的后验；形式仍是“模型输出减后验”，只是多了格点占用作为权重，模型很少经过的格点几乎不受更新。在上例的格点$(1,1)$，占用为$.6\times.675/.6066=.668$；空白边后验为$.6\times.5\times.72/.6066=.356$，标签边后验为$.6\times.5\times.63/.6066=.312$。空白的梯度为$.668\times.5-.356=-.022$，$a$的梯度为$+.022$，更新会略微提高这一格点的空白概率，因为经空白到达的$(2,1)$有更大的后续概率。

推断时，贪心解码在当前格点取概率最大的符号：空白使$t$加1，非空白追加到前缀、更新预测网络状态并保持$t$不变。束搜索需要处理同一前缀可经不同路径到达的情况。算法\ref{alg:nlp-speech-rnnt-beam}是Graves原文算法的简化形式，省略了原文中把短前缀的概率转移给长前缀的合并步骤。

\begin{algorithm}[RNN-T束搜索（简化形式）]
  \label{alg:nlp-speech-rnnt-beam}
  \begin{algorithmic}[1]
    \Require 编码输出$\bm{h}_1,\ldots,\bm{h}_T$，预测网络与联合网络，束宽$K$
    \Ensure 识别文本
    \State $B\gets\{\epsilon\}$，$P(\epsilon)=1$，缓存预测网络输出$\bm{g}(\epsilon)$
    \For{$t=1,\ldots,T$}
      \State $A\gets B$，再把$B$置空
      \While{$A$中最大概率高于$B$中第$K$大的概率（$B$不足$K$项时视为0）}
        \State 从$A$中取出概率最大的前缀$y^*$
        \State 由$\bm{h}_t$与缓存的$\bm{g}(y^*)$按式\eqref{eq:nlp-speech-rnnt-joint}计算分布$q$
        \State 把$y^*$以概率$P(y^*)q(\varnothing)$加入$B$，已存在则概率相加 \Comment{空白：进入下一步}
        \For{$k\in V$}
          \State 把$y^*k$以概率$P(y^*)q(k)$加入$A$，已存在则相加；计算并缓存$\bm{g}(y^*k)$
        \EndFor
      \EndWhile
      \State 保留$B$中概率最大的$K$个前缀
    \EndFor
    \State \Return $B$中$\log P(y)/|y|$最大的$y$
  \end{algorithmic}
\end{algorithm}

每次扩展都乘以小于1的概率，$A$中的最大概率不断下降，循环必然结束。实际系统通常还限制每个声学步最多输出的符号数，并只对概率最高的若干个$k$扩展；达到上限后强制进入下一步，会改变搜索结果，这是实现上的折中。预测网络的状态只取决于前缀，所以相同前缀共享一份缓存；预测网络很小，这部分计算远少于编码器。

\subsubsection{显存开销与剪枝训练}
\label{sec:nlp-speech-rnnt-pruned}

RNN-T训练的主要瓶颈是联合网络的四维输出。取批大小$B=16$，10秒声音经4倍降采样后$T=250$，目标长$U=60$，子词词表500个加空白，输出张量有$16\times250\times61\times501\approx1.22\times10^8$个元素，32位浮点约0.49\,GB；反向传播还要保存同样大小的梯度，以及$d_{\mathrm{joint}}=512$时约同样规模的隐藏张量。若改用约5000个汉字的词表，仅输出张量就约4.9\,GB。显存随$T$、$U$与$|V|$的乘积增长，这是CTC与AED都没有的代价。

一类做法从实现入手。函数合并把softmax、损失与梯度合在一个核函数里计算，梯度直接覆盖原来的输出缓冲，不再单独保存概率张量\scite{nlp-speech-li2019rnnt}；按长度排序组批则减少填充造成的浪费。这些做法降低常数，但不改变$T\times U\times|V|$的规模。

\term{剪枝RNN-T}（pruned RNN-T）从格点本身入手\scite{nlp-speech-kuang2022}。在整个格点中，大多数$(t,u)$的占用接近0：第$t$个声学步附近实际可能输出的只是目标中少数几个位置。它分两步计算。第一步用一个简化的加性联合网络：编码器与预测网络各自投影到词表维度，得到$\bm{l}^{\mathrm{enc}}_t$与$\bm{l}^{\mathrm{pred}}_u$，令
\begin{equation}
 \tilde{z}_{t,u,k}=l^{\mathrm{enc}}_{t,k}+l^{\mathrm{pred}}_{u,k}.
 \label{eq:nlp-speech-pruned-simple}
\end{equation}
格点递推只需要空白和目标符号这两个分数，以及softmax的归一化项$\log\sum_k\exp(l^{\mathrm{enc}}_{t,k}+l^{\mathrm{pred}}_{u,k})$；后者可以写成两个矩阵$\exp(\bm{L}^{\mathrm{enc}})\in\mathbb{R}^{T\times|V|}$与$\exp(\bm{L}^{\mathrm{pred}})^{\top}$的乘积，不必展开$T\times U\times|V|$的张量。用这个简化格点算出损失及其对各格点的梯度，就知道每个$t$上哪些$u$的占用较大。第二步为每个$t$选一个起点$p_t$，要求$p_t$随$t$不减，以保持单调；只在宽度为$S$的带$u\in[p_t,p_t+S)$内计算式\eqref{eq:nlp-speech-rnnt-joint}的完整联合网络及其损失。总损失是剪枝损失与简化损失的加权和，原文发现把简化损失乘以0.5效果最好，它同时训练用于选带的投影。

带宽$S$通常取个位数。上例取$S=5$、汉字词表时，输出张量从约4.9\,GB降到$16\times250\times5\times5001\times4$字节，约0.40\,GB。代价是选带依赖简化模型的质量：训练初期简化模型尚不可靠，带选得不准会使真正占用大的格点落在带外。开源实现icefall因此在最初几千步只用简化损失，随后以较小权重加入剪枝损失，再升到完整权重；带宽过窄也会影响收敛。

RNN-T同时保留了单调路径求和与输出历史，编码器做成流式后就能逐步输出，是目前在线识别的主流结构之一。它的代价集中在训练显存和实现复杂度上；预测网络只读文字，外部语言模型融合时还要考虑它已经学到的内部语言先验（第\ref{sec:nlp-speech-fusion}小节）。

\subsection{大规模弱监督识别：Whisper}
\label{sec:nlp-speech-whisper}

前述模型通常在一两个人工转写的数据集上训练。它们在同分布测试集上错误率很低，换到另一种口音、噪声或领域时却明显变差，而人工转写难以扩大到覆盖这些条件。Whisper换了一个方向：结构仍是标准的Transformer编码器—解码器，训练数据改为约68万小时从网络收集的音频与字幕配对\scite{radford2023whisper}。其中约11.7万小时覆盖英语以外的96种语言，约12.5万小时是其他语言音频配英语译文。这些字幕没有经过人工核对，属于弱监督；数据处理用启发式规则删去疑似机器生成的转写（例如全大写或全无标点的文本），并删去字幕语种与音频语种不符的样本。

模型输入是30秒窗口。音频重采样到16\,kHz，计算25\,ms窗长、10\,ms帧移的80维log-Mel，不足30秒的补零，得到$\bm{X}\in\mathbb{R}^{3000\times80}$。卷积前端由两层宽度为3的一维卷积组成，激活函数为GELU：第一层步长1，输出$3000\times d$；第二层步长2，输出$1500\times d$，每个位置对应20\,ms。加上正弦位置编码后进入预归一化的Transformer编码器，输出$1500\times d$。解码器使用可学习的位置编码，输入与输出词嵌入共享参数，每层通过交叉注意力（式\eqref{eq:nlp-llm-cross-attention}）读取编码器输出。每层的自注意力、前馈与残差结构与第\ref{sec:nlp-llm-decoder-layer}小节的解码器层相同，只是编码器不用因果掩码。模型从约3900万参数（4层，$d=384$）到约15.5亿参数（32层，$d=1280$）分为多档。与第\ref{sec:nlp-speech-encoder}小节Conformer常用的4倍二维卷积降采样相比，Whisper只降采样2倍，编码序列更长，但30秒的固定窗口限制了总长度。

同一段声音可以要求转写、译成英语或只判断语种，Whisper把任务写成解码器开头的特殊词元。一个带时间戳的中文转写目标形如
\begin{center}
\small
\texttt{<|startofprev|>} 机器学习读书会通知 \texttt{<|startoftranscript|><|zh|><|transcribe|>}\\
\texttt{<|0.00|>} 活动改到四月十八日 \texttt{<|2.40|><|2.40|>} 下午两点在海棠楼 \texttt{<|4.60|><|endoftext|>}
\end{center}
\texttt{<|startofprev|>}之后是前一窗口的转写，作为上下文；\texttt{<|zh|>}是语种，\texttt{<|transcribe|>}与\texttt{<|translate|>}区分转写和译成英语；每个文本片段前后各有一个时间词元，以20\,ms为间隔量化窗口内的相对时间。不需要时间时改用\texttt{<|notimestamps|>}；窗口内没有语音时，语种位置预测\texttt{<|nospeech|>}。训练损失是目标词元的条件负对数概率，与式\eqref{eq:nlp-speech-las}相同，但上下文部分不计损失。这样，语种识别、有无语音判断、转写、翻译和时间戳都成为同一个解码器的生成任务。

长录音按30秒窗口依次处理：用当前窗口最后一个完整片段的结束时间决定下一窗口的起点，把已解码文字作为下一窗口的上下文，局部时间戳加上窗口在录音中的偏移。原文还采用几项启发式：束宽为5的束搜索；输出的平均对数概率过低或文本压缩率过高（提示重复）时，提高采样温度重新解码；\texttt{<|nospeech|>}概率高且平均对数概率低时，把该窗口判为无语音。这些规则说明，固定窗口的自回归解码在长音频上需要额外的保护。

大规模数据带来的主要收益是跨领域的鲁棒性。原文比较了在LibriSpeech上错误率相近的模型：Whisper不在任何目标数据集上微调，在其他多个英语识别数据集上的平均相对错误率比在LibriSpeech上有监督训练的模型低约55\%。代价同样明显。第一，解码器本质上是以声音为条件的语言模型，声音证据弱时语言先验可能占上风：静音或背景音乐段可能生成“谢谢观看”一类网络字幕中常见的句子，或在窗口之间重复上一段文字\scite{nlp-speech-koenecke2024}。CTC与RNN-T的每个非空白输出都要在某个声学步上与空白竞争，没有声音支持时默认输出空白，这类无依据的长段插入少得多，但并非不可能。第二，固定30秒窗口、全局注意力和自回归解码都要求先拿到整段声音，Whisper不能直接流式运行，需要第\ref{sec:nlp-speech-whisper-streaming}小节的外部组织。第三，把任务词元设错，例如把转写设为翻译，输出语言就会改变，这不是听错了声音，评价时应与声学错误分开统计。

\subsection{流式识别}
\label{sec:nlp-speech-streaming}

在线识别要求每个单位在其声音结束后不久提交，并且已提交的结果不再频繁改动。延迟有两个来源：模型为了看见右侧声音而等待，以及计算本身所需的时间。本小节分两条路线。第一条从模型结构入手，限制编码器的右侧可见范围，使CTC或RNN-T能逐块输出。第二条处理Whisper这类整段模型，反复重新识别不断增长的缓冲，只提交多次假设一致的前缀。两者都要同时报告提交延迟、修订量和最终错误率，因为结果稳定不等于结果正确。

\subsubsection{分块注意力与流式编码器}
\label{sec:nlp-speech-chunk-attention}

自注意力编码器的每个位置默认读取整段序列，最直接的流式改造是用注意力掩码限制可见范围。若每层允许看右侧$r$步，$N$层叠加后最上层的位置实际依赖右侧$Nr$步：12层、每层2步、每步40\,ms时，前瞻达到960\,ms，延迟随层数累积。\term{分块注意力}（chunk attention）改为把编码序列切成长为$C$的块，每个位置可以看本块全部位置和左侧$L$个历史块。以从1开始的位置索引，记$c(i)=\lfloor(i-1)/C\rfloor$为位置$i$所在的块号，掩码为
\begin{equation}
 M_{i,j}=\begin{cases}
   0, & c(i)-L\leq c(j)\leq c(i),\\
   -\infty, & \text{其他},
 \end{cases}
 \qquad
 \operatorname{Attn}=\operatorname{softmax}\!\left(\frac{\bm{Q}\bm{K}^{\top}}{\sqrt{d_k}}+\bm{M}\right)\bm{V}.
 \label{eq:nlp-speech-chunk-mask}
\end{equation}
图\ref{fig:nlp-speech-chunk-mask}画出了$C=2$、$L=1$时的掩码。由于每层的可见范围都止于本块末尾，叠加多少层，右侧都不会越过块边界。块内第一个位置要等到块末尾才能计算，最多等待$C-1$步，块大小因而直接决定了模型带来的延迟下界：$C=16$个40\,ms编码步时，最长等待约640\,ms，与层数无关。

\begin{figure}[htbp]
  \centering
  % 第九卷 15.2.7.1 节：分块注意力掩码（正文栏宽，约60 mm）。
% 8个编码步，块大小C=2，左侧保留L=1个历史块。蓝色格为可见的键位置。
\begin{tikzpicture}[x=1mm,y=1mm,figure picture]
  \foreach \i in {0,...,7} {
    \foreach \j in {0,...,7} {
      \pgfmathtruncatemacro{\ci}{floor(\i/2)}
      \pgfmathtruncatemacro{\cj}{floor(\j/2)}
      \pgfmathtruncatemacro{\vis}{and(\cj<=\ci,\cj>=\ci-1)}
      \ifnum\vis=1
        \path[fill=FigureInputFill,draw=FigureInput,line width=.4pt]
          ({5*\j},{-5*\i}) rectangle ++(5,-5);
      \else
        \path[fill=white,draw=FigureGrid,line width=.4pt]
          ({5*\j},{-5*\i}) rectangle ++(5,-5);
      \fi
    }
    \pgfmathtruncatemacro{\k}{\i+1}
    \node[font=\figurefont\scriptsize,anchor=east] at (-1,{-5*\i-2.5}) {$\k$};
    \node[font=\figurefont\scriptsize,anchor=south] at ({5*\i+2.5},1) {$\k$};
  }
  \foreach \k in {0,2,4,6,8} {
    \draw[FigureStroke,line width=.8pt] ({5*\k},0) -- ({5*\k},-40);
    \draw[FigureStroke,line width=.8pt] (0,{-5*\k}) -- (40,{-5*\k});
  }
  \node[font=\figurefont\scriptsize,anchor=south] at (20,6) {键位置$j$};
  \node[font=\figurefont\scriptsize,rotate=90,anchor=south] at (-6,-20) {查询位置$i$};
  \node[font=\figurefont\scriptsize,anchor=west,align=left,text=FigureMuted]
    at (43,-20) {块大小$C=2$\\左侧历史$L=1$块\\粗线为块边界};
\end{tikzpicture}

  \caption{分块注意力掩码，8个编码步，块大小$C=2$，左侧保留$L=1$个历史块。第$i$行的蓝色格表示查询位置$i$可以读取的键位置。每个位置都能看到本块内右侧的位置，但不能越过块边界；各层使用同一掩码，叠加后的可见范围仍止于块末尾。}
  \label{fig:nlp-speech-chunk-mask}
\end{figure}

块大小是准确率与延迟的权衡，而不同应用的延迟预算不同。U2的\term{动态块训练}（dynamic chunk training）为每个批次随机选择块大小：以一半概率使用整句（全上下文），否则从1至25个编码步中均匀抽取，使同一个模型在推断时可以按延迟预算选择块大小\scite{nlp-speech-zhang2020u2}。该工作还把它与第\ref{sec:nlp-speech-joint-ctc-attention}小节的联合模型结合成两遍解码：CTC分支在分块编码器上流式输出暂定结果，整句结束后由注意力解码器对CTC的$n$个最佳候选重打分，作为最终结果。

Conformer中的其他模块也要随之改造。卷积模块的深度卷积若两侧补零，会读到右侧未来的帧，需改为只在左侧补零的因果卷积；跨块时缓存上一块末尾的$k-1$帧（$k$为卷积核宽度），作为本块卷积的左侧输入。注意力同样缓存左侧历史块的键与值，与第\ref{sec:nlp-inference-cache}小节文本生成中的键值缓存作用相同。批归一化在推断时使用训练中累积的统计量，本身不读未来帧，但分块训练与整句训练的统计差异可能造成不一致，一些实现改用层归一化。CTC与RNN-T接上这种编码器即可逐块输出；AED的全局注意力解码不能直接这样改造，这是流式场景偏向CTC与RNN-T的主要原因。

\subsubsection{稳定前缀提交：Whisper-Streaming}
\label{sec:nlp-speech-whisper-streaming}

整段识别器也可以通过缓冲和反复识别接入实时声音，难点是决定哪些结果已经可以提交。Whisper-Streaming采用\term{局部一致性}（LocalAgreement）策略：对不断增长的声音缓冲反复识别，比较连续两次的结果，只提交它们一致的公共前缀，其余部分留待更多声音到来后再确认\scite{nlp-speech-machacek2023}。这是模型之外的调度规则，Whisper本身的概率模型并不包含它；其思路与第\ref{sec:nlp-generation-adaptive-policy}小节同声翻译中的自适应读写策略相似。

每轮收集到至少规定长度的新声音后，识别当前缓冲；跳过已经提交的文字，对上次与本次的剩余部分求最长公共前缀，扩展提交边界。已提交部分仍留在缓冲内为识别提供上下文，但它在新识别中的变化不再对外撤回。缓冲过长时，在已提交的句末位置裁掉旧声音，并把已提交文字作为下一轮的提示。若一轮计算慢于新声音的积累，下一轮会处理更多声音，实际等待由输入速度与计算速度共同决定。

设三次教学假设依次为“改到／四月／十七日”“改到／四月／十八日”“改到／四月／十八日／下午”。第一次没有可比较的结果；第二次提交“改到／四月”，日期两次不同而保留；第三次才提交“十八日”。若第一次的“十七日”被直接播放给听者，后来修改屏幕文字也无法撤销已经听到的日期。反过来，若连续两次都误识为“十七日”，局部一致性同样会提交错误，它判断的是稳定性而不是正确性。评价时固定块长、缓冲裁剪和计分规范，分别统计最终字符错误率、暂定文字的修订量和各单位的提交延迟；被漏识的字没有提交时刻，须同时报告遗漏，不能只平均已识别字的延迟。

\subsection{识别模型的比较}
\label{sec:nlp-speech-asr-comparison}

表\ref{tab:nlp-speech-asr-comparison}汇总了本节的六类识别模型。各列依次回答：训练需要什么标签；输出之间是否假设条件独立；对齐是否被约束为单调；能否在声音到达时逐步输出；外部语言模型怎样接入；能否直接得到每个单位的时间位置；主要的计算或工程代价是什么。

\begin{table*}[htbp]
  \centering
  \footnotesize
  \begin{tabularx}{\linewidth}{@{}>{\raggedright\arraybackslash}p{17mm}>{\raggedright\arraybackslash}p{16mm}>{\raggedright\arraybackslash}p{13mm}>{\raggedright\arraybackslash}p{14mm}>{\raggedright\arraybackslash}p{17mm}>{\raggedright\arraybackslash}p{19mm}>{\raggedright\arraybackslash}p{17mm}>{\raggedright\arraybackslash}X@{}}
    \toprule
    方法 & 训练标签 & 输出条件独立 & 单调对齐 & 流式 & 外部语言模型 & 对齐与时间 & 主要代价\\
    \midrule
    混合HMM--DNN & 帧级状态对齐（来自HMM--GMM） & 是，给定状态各帧独立 & 是，由HMM拓扑保证 & 可以，受网络输入窗口限制 & 解码图中的$G$，可词格重打分 & 状态级对齐直接可得 & 多阶段训练，需发音词典\\
    CTC & 整句文字 & 是 & 是 & 配合分块或因果编码器 & 前缀束搜索中按字加分 & 尖峰位置，偏离发音边界 & 需外部语言模型补历史\\
    AED（LAS） & 整句文字 & 否 & 否，注意力可跳读 & 难，需整段编码 & 束搜索中浅层融合 & 注意力权重，不可靠 & 跳读、重复与过早结束\\
    联合CTC/\allowbreak 注意力 & 整句文字 & 注意力分支否 & CTC分支约束 & 两遍解码时首遍可流式 & 联合分数中加一项 & 可由CTC分支取得 & 两个分支都要计算\\
    RNN-T & 整句文字 & 否 & 是 & 配合分块或因果编码器 & 浅层融合，宜扣除内部语言先验 & 发出位置，常有延迟 & $B\times T\times U\times|V|$显存\\
    Whisper & 弱监督字幕，约68万小时 & 否 & 否 & 不能直接流式，需外部调度 & 可融合，但解码器已含强语言先验 & 20\,ms时间词元，片段级 & 无语音段幻觉，30秒固定窗口\\
    \bottomrule
  \end{tabularx}
  \caption{六类语音识别模型的对照。“流式”指在声音到达过程中逐步给出稳定输出的能力，取决于编码器可见范围与解码方式，不由损失函数单独决定；“对齐与时间”指不另做强制对齐时能直接取得的时间信息。}
  \label{tab:nlp-speech-asr-comparison}
\end{table*}

按表中的列可以读出各方法依次解决了什么问题。CTC用路径求和去掉了帧级对齐标签与多阶段训练；AED引入输出历史，却丢掉了单调约束；联合训练用CTC分支把单调约束找回来；RNN-T在一个模型里同时得到输出历史、单调对齐和流式能力，代价转移到训练显存；Whisper不改变结构，而是用弱监督数据扩大覆盖，换来鲁棒性，同时带来幻觉与非流式的问题。选型取决于延迟预算、可用数据与部署资源：需要低延迟在线输出时，分块编码器上的RNN-T或CTC更合适；离线转写、语种多样且没有领域内标注时，大规模弱监督模型更合适；需要精确时间边界或只有少量数据时，HMM路线仍有用处。没有一种方法在所有条件下都占优。
